\documentclass[11pt,a4paper]{article}

\usepackage[utf8]{inputenc}
\usepackage[T1]{fontenc}
\usepackage[brazilian]{babel}
\usepackage[margin=2cm]{geometry}
\usepackage{enumitem}
\usepackage{hyperref}
\usepackage{titlesec}
\usepackage{parskip}

\hypersetup{
  colorlinks=true,
  linkcolor=black,
  urlcolor=black,
  pdftitle={Kamille Hartmann Maccagnan - Currículo},
  pdfauthor={Kamille Hartmann Maccagnan}
}

\titleformat{\section}{\large\bfseries}{}{0em}{}[\titlerule]
\titlespacing*{\section}{0pt}{12pt}{6pt}

\pagestyle{empty}
\setlength{\parindent}{0pt}
\setlength{\parskip}{6pt}

\newcommand{\cvheader}[4]{%
  \begin{center}
    {\LARGE\bfseries #1}\\[4pt]
    {\large #2}\\[6pt]
    #3\\[2pt]
    #4
  \end{center}
  \vspace{4pt}
}

\newcommand{\experience}[3]{%
  \textbf{#1} \hfill \textit{#2}\\[4pt]
  #3
  \vspace{8pt}
}

\begin{document}

\cvheader{Kamille Hartmann Maccagnan}%
{Analista de Projetos Jr | PMO | Indicadores e Dashboards | Governança | Serviços de TI}%
{Campo Comprido, Curitiba-PR \textbullet{} (41) 9 9928-4424 \textbullet{} \href{mailto:kamillehartmannm@gmail.com}{kamillehartmannm@gmail.com}}%
{\href{https://linkedin.com/in/kamille-hartmann-maccagnan/}{linkedin.com/in/kamille-hartmann-maccagnan}}

\section*{Resumo Profissional}

Profissional organizada e analítica, com experiência em PMO, acompanhamento de projetos e apoio à governança de entregas em ambiente corporativo e multinacional. Atuação na consolidação de informações, elaboração de dashboards e relatórios de acompanhamento, monitoramento de indicadores, cronogramas, riscos e pendências, com comunicação clara entre equipes técnicas, gestores e clientes. Vivência em documentação operacional, procedimentos, bases de conhecimento e condução de transferência de conhecimento por meio de treinamentos e materiais de apoio. Uso diário de SQL para consultas e extração de dados, Power BI para visualização de indicadores e ferramentas de Inteligência Artificial, como ChatGPT e Copilot, para organização de informações, elaboração de documentos e geração de resumos executivos. Perfil colaborativo, proativo e orientado a resultados, com atenção aos detalhes, facilidade para transformar informações complexas em apresentações objetivas e interesse em Gestão de Projetos e Serviços de TI.

\section*{Experiência Profissional}

\experience{Anjuss -- Analista de Suporte de TI}{07/2025 -- Atual}{%
Atuação no suporte a projetos e operações de T.I. em ambiente de varejo, com consolidação de demandas, organização de pendências e acompanhamento de entregas entre lojas, operações e tecnologia. Elaboração e manutenção de dashboards em Power BI e relatórios recorrentes para monitoramento de indicadores, visibilidade de status e apoio à tomada de decisão. Responsável pela documentação operacional estruturada, incluindo POPs, registros de atividades e bases de conhecimento para repasse entre equipes. Condução de treinamentos, onboarding e transferência de conhecimento para novos funcionários e área de T.I., promovendo continuidade operacional e autonomia no uso dos sistemas. Participação em alinhamentos com usuários via WhatsApp corporativo e reuniões de acompanhamento, com comunicação clara de solicitações, prazos e retornos. Uso diário de SQL para consultas e extração de informações e de ferramentas de Inteligência Artificial para organização de demandas, elaboração de comunicações e resumos de acompanhamento. Apoio à investigação de incidentes, priorização de planos de ação e escalonamento para times técnicos ou fornecedores quando necessário.}

\experience{HORSE do Brasil -- Estagiária de TI (IS/IT LATAM)}{07/2024 -- 07/2025}{%
Atuação em PMO e apoio à governança de projetos em ambiente multinacional, com acompanhamento de cronogramas, marcos de entrega, indicadores de desempenho, riscos e pendências de múltiplas iniciativas. Consolidação de informações provenientes de diferentes frentes e elaboração de relatórios executivos e materiais de acompanhamento para apresentação objetiva a lideranças, clientes internos, fornecedores e times técnicos. Participação em reuniões de status e alinhamento conduzidas em até três idiomas simultaneamente, facilitando comunicação entre áreas diversas e apoio à execução eficiente das entregas. Atuação na documentação de processos, análise de desvios operacionais, identificação de impedimentos e evolução de práticas de gestão com foco em qualidade e melhoria contínua.}

\experience{Restaurante Familiar -- Analista de Dados e Suporte Técnico}{01/2023 -- 07/2024}{%
Atuação no apoio à gestão do negócio com consolidação de informações operacionais e financeiras, estruturação de base de dados em Excel e dashboards em Power BI para acompanhamento de indicadores e relatórios de performance. Responsável pela padronização de informações, identificação de inconsistências, documentação de procedimentos e apresentação de análises de forma clara para apoiar decisões entre áreas operacionais e administrativas.}

\section*{Projetos e Conquistas}

Estruturação de KPIs e dashboards para monitoramento de indicadores e suporte à governança de entregas. Participação em iniciativas de melhoria contínua e otimização de processos. 1º lugar no Transformation Day Renault 2023 com solução orientada a dados e melhoria de processos. Participação no Hackathon Bosch 2023 com foco em resolução de problemas de negócio com uso de dados.

\section*{Habilidades Técnicas}

\textbf{Gestão de Projetos e PMO:} Acompanhamento de cronogramas, marcos, riscos, pendências e planos de ação, consolidação de informações, governança de projetos, reuniões de status, Scrum, Kanban, Boas Práticas PMOBOK.

\textbf{Indicadores, Dashboards e Relatórios:} Power BI (DAX, Power Query, insights, Q\&A), KPIs, relatórios executivos e de acompanhamento, SQL (consultas e extração diária), Excel, transformação de informações complexas em visões objetivas.

\textbf{Documentação e Transferência de Conhecimento:} POPs, procedimentos operacionais, bases de conhecimento, treinamentos, onboarding, repasse de conhecimento entre equipes, Notion.

\textbf{Serviços de TI e Operações:} Suporte técnico, Service Desk, registro e rastreabilidade de demandas, noções de ITIL e gestão de serviços, ServiceNow.

\textbf{Ferramentas de Gestão:} Monday, Trello, Notion, Spotfire, Pacote Office.

\textbf{Inteligência Artificial:} ChatGPT, Copilot, apoio à produtividade, organização de informações, elaboração de documentos e resumos executivos.

\textbf{Idiomas:} Inglês avançado \textbullet{} Espanhol básico \textbullet{} Coreano básico.

\section*{Capacitações Complementares}

Capacitação em aproveitamento de ferramentas de Inteligência Artificial.

\section*{Formação Acadêmica}

\textbf{PUCPR} -- Graduação em Gestão da Tecnologia da Informação \hfill Previsão de conclusão: 06/2026\\
Curitiba, PR

\textbf{Escola Conquer} -- Pós-graduação em Liderança e Gestão em Tecnologia \hfill Conclusão: 2024\\
Curitiba, PR

\textbf{Universidade Tuiuti do Paraná} -- Graduação em Medicina Veterinária \hfill Conclusão: 06/2019\\
Curitiba, PR

\textbf{Qualittas} -- Pós-graduação em Clínica Médica e Cirúrgica de Animais Exóticos e Pets Não Convencionais \hfill Conclusão: 12/2022\\
Curitiba, PR

\end{document}
